% Source: source/普通高考/2011/2011新课标文(河南,黑龙江,吉林,海南,宁夏,山西,新疆).pdf,
% page 1, question 5. pdfimages -list reports no embedded bitmap; the chart is vector/text artwork.
% Grid evidence: tmp/review_flowchart_2011_new_curriculum_liberal/grid/page-1-grid100.png
% (100px coarse + 50px fine) and q5-block-grid50.png (50px coarse + 25px fine).
% Comparison crop: tmp/review_flowchart_2011_new_curriculum_liberal/crop/q5_original_final.png,
% taken from the ungridded page render with a safety border.
% Nodes: 开始; 输入 N; k=1,p=1; p=p·k; k<N?; k=k+1; 输出 p; 结束.
% Directed edges: 开始→输入→初始化→p更新→判定; 判定(是)→k更新→右上回线→p更新;
% 判定(否)→输出 p→结束. The loop return is separate from the output branch.
\begin{tikzpicture}[
  >=Stealth,
  line width=0.55pt,
  line cap=round,
  line join=round,
  every node/.style={font=\normalsize,anchor=center,align=center,
    execute at begin node={\setlength{\parindent}{0pt}\noindent}},
  flowbox/.style={draw,minimum height=.68cm,inner xsep=4pt,inner ysep=2pt,
    anchor=center,align=center,
    execute at begin node={\setlength{\parindent}{0pt}\noindent}},
  terminal/.style={flowbox,rounded corners=.18cm,minimum width=1.35cm,
    minimum height=.72cm},
  io/.style={flowbox,shape=trapezium,trapezium left angle=72,
    trapezium right angle=108,minimum height=.78cm},
  decision/.style={flowbox,shape=diamond,aspect=2.75,inner sep=2pt},
  flowarrow/.style={->,line width=0.55pt},
]
  \node[terminal] (start) at (0,9.25) {开始};
  \node[io,minimum width=2.35cm] (input) at (0,7.92) {输入 $N$};
  \node[flowbox,minimum width=3.30cm] (init) at (0,6.62)
    {$k=1,p=1$};
  \coordinate (loopjoin) at (0,5.76);
  \node[flowbox,minimum width=2.80cm] (product) at (0,4.82)
    {$p=p\cdot k$};
  \node[decision,minimum width=3.40cm,minimum height=.98cm] (test) at (0,3.38)
    {$k<N$};
  \node[flowbox,minimum width=2.80cm] (increment) at (3.20,4.22)
    {$k=k+1$};
  \node[io,minimum width=2.05cm] (output) at (0,1.92) {输出 $p$};
  \node[terminal] (finish) at (0,0.62) {结束};

  \draw[flowarrow] (start.south) -- (input.north);
  \draw[flowarrow] (input.south) -- (init.north);
  \draw (init.south) -- (loopjoin);
  \draw[flowarrow] (loopjoin) -- (product.north);
  \draw[flowarrow] (product.south) -- (test.north);
  \draw[flowarrow] (test.south) -- node[right=2pt] {否} (output.north);
  \draw[flowarrow] (output.south) -- (finish.north);

  % Affirmative branch updates k and returns along the upper right route.
  \coordinate (right-branch) at (3.20,3.38);
  \coordinate (right-top) at (3.20,5.76);
  \coordinate (loop-tip) at (-0.60,5.76);
  \draw (test.east) -- node[below=2pt,inner sep=0pt] {是} (right-branch);
  \draw[flowarrow] (right-branch) -- (increment.south);
  \draw (increment.north) -- (right-top);
  \draw[flowarrow] (right-top) -- (loop-tip) -- (loopjoin);
  \path[use as bounding box] (-2.55,0.30) rectangle (4.55,9.60);
\end{tikzpicture}
